% Complete Section 3.4 and Box 3.1, original PDF 37-42 / printed 27-32.
% The floating summary is placed after the final main-text paragraph.
% block: vector-operators-heading page:37 type:section
\subsection{{Example: grad, curl 和 div 及其伴随}}\label{sec:vector-operators-adjoints}

% block: vector-operators-intro page:37 type:ordinary-paragraph
{本节考虑 de Rham 复形中的三个向量微积分算子, 说明如何将它们视为 Hilbert 空间上的闭无界算子, 并在这一框架下推导其基本性质.}

% block: distributional-gradient-domain page:37 type:ordinary-paragraph
{设 $\Omega$ 表示 $\mathbb R^n$ 中的有界区域. 先考虑从 $L^2(\Omega;\mathbb R^n)$ 到 $L^2(\Omega)$ 的无界算子, 其定义域为 $C_0^\infty(\Omega;\mathbb R^n)$, 对 $\phi\in C_0^\infty(\Omega;\mathbb R^n)$, 由 $\phi\mapsto-\operatorname{div}\phi$ 给出. 为明确其定义域, 将此算子记为 $(-\operatorname{div},C_0^\infty)$. 尽管它不是闭算子, 但显然稠密定义, 故其伴随 $(-\operatorname{div},C_0^\infty)^*$ 是闭算子 $L^2(\Omega)\to L^2(\Omega;\mathbb R^n)$. 按定义, $\Omega$ 上的实值函数 $u$ 属于该伴随的定义域, 是指它属于 $L^2(\Omega)$, 且存在 $L^2(\Omega;\mathbb R^n)$ 中的向量场 $v$, 使得}
\[
-\int_\Omega u\operatorname{div}\phi\,dx=\int_\Omega v\cdot\phi\,dx,\qquad\phi\in C_0^\infty(\Omega;\mathbb R^n).
\]

% block: gradient-hilbert-space page:37 type:ordinary-paragraph
{而此等式恰好表明 $u$ 的分布梯度 $\operatorname{grad}u$ 属于 $L^2$ (即 $\operatorname{grad}u=v=(-\operatorname{div},C_0^\infty)^*u$). 因此, $(-\operatorname{div},C_0^\infty)^*$ 的定义域是所有分布梯度属于 $L^2$ 的 $L^2$ 函数的集合, 这恰好就是 Sobolev 空间 $H^1(\Omega)$ 的定义. 我们已经证明}
\[
(-\operatorname{div},C_0^\infty)^*=(\operatorname{grad},H^1).
\]
{算子 $(\operatorname{grad},H^1)$ 不仅是闭的, 而且显然也稠密定义 (例如其定义域包含 $C_0^\infty$, 后者已在 $L^2$ 中稠密). 它是闭算子这一事实确立了 $H^1$ 的完备性.}

% block: divergence-hilbert-space page:37 type:ordinary-paragraph
{现交换 div 和 grad 的角色, 考虑从 $L^2(\Omega)$ 到 $L^2(\Omega;\mathbb R^n)$ 的无界算子 $(-\operatorname{grad},C_0^\infty)$. 与上面完全相同地推理, 有}
\[
(-\operatorname{grad},C_0^\infty)^*=(\operatorname{div},H(\operatorname{div})),
\]
{其中右端算子是闭的且稠密定义. 其定义域为 Hilbert 空间 $H(\operatorname{div};\Omega)$, 由分布散度属于 $L^2$ 的 $L^2$ 向量场组成.}

% block: grad-div-integration-by-parts-intro page:38 type:ordinary-paragraph
{下一个目标是计算刚才得到的两个算子 $(\operatorname{grad},H^1)$ 和 $(\operatorname{div},H(\operatorname{div}))$ 各自的伴随. 关键在于分部积分公式}
\begin{equation}\label{eq:grad-div-classical-green}
\int u\operatorname{div}w\,dx+\int\operatorname{grad}u\cdot w\,dx=\int_{\partial\Omega}u|_{\partial\Omega}w|_{\partial\Omega}\cdot n\,ds.
\end{equation}
{当 $u\in C^1(\overline\Omega)$, $w\in C^1(\overline\Omega;\mathbb R^n)$ 且边界光滑时, 此式当然成立; 但当 $w$ 仅属于 $H(\operatorname{div})$ 时, $w$ 在 $\partial\Omega$ 上的迹这类量是否有意义尚不清楚. 因此先研究这一问题. 关键结论见 Theorem~\ref{thm:normal-trace-extension}.}

% block: sobolev-trace-facts page:38 type:ordinary-paragraph
{为此, 并在本节余下部分, 假设 $\Omega$ 具有 Lipschitz 边界. 先回顾 Sobolev 空间 $H^1$ 的几个基本事实. 第一, 空间 $C^\infty(\overline\Omega)$ 在 $H^1(\Omega)$ 中稠密 [\BookNumber{41}, Section 5.3, Theorem 2]. 第二, 限制算子 $C^1(\overline\Omega)\to L^2(\partial\Omega)$ 可延拓为有界迹算子 $\operatorname{tr}:H^1(\Omega)\to L^2(\partial\Omega)$ [\BookNumber{41}, Section 5.5, Theorem 1]. 迹算子的像空间为 $H^{1/2}(\partial\Omega)$; 事实上, 这也是定义该空间的一种方式. 按这种方式, $H^{1/2}(\partial\Omega)$ 中的范数由下式给出:}
\begin{equation}\label{eq:trace-quotient-norm}
\|g\|_{1/2,\partial\Omega}=\inf_{\operatorname{tr}u=g}\|u\|_{1,\Omega}.
\end{equation}

% block: trace-annihilator-intro page:38 type:ordinary-paragraph
{第三, 迹算子的核为空间 $\mathring H^1(\Omega)$, 它定义为 $C_0^\infty(\Omega)$ 在 $H^1(\Omega)$ 中的闭包. 见 [\BookNumber{41}, Section 5.5, Theorem 2]. 已知作用于 $H^1(\Omega)$ 的迹算子的像与核后, 便可利用对偶刻画那些在核 $\mathring H^1(\Omega)$ 上为零的 $H^1(\Omega)$ 上的泛函. 下述结论是初等的.}

% block: dual-surjection-annihilator page:38 type:lemma
\Needspace{6\baselineskip}
\begin{Lemma}\label{lem:dual-surjection-annihilator}
{若 $\gamma:X\to Y$ 是赋范线性空间之间的有界线性满射, 核为 $Z$, 则对偶映射 $\gamma':Y'\to X'$ 是有界单射, 其像等于 $Z$ 的零化子 (定义为 $\{\,L\in X'\mid L|_Z\equiv0\,\}$). 因此, 对零化子中的 $L$, 存在唯一 $g\in Y'$, 使得}
\[
Lu=\langle g,\gamma u\rangle,\qquad u\in X.
\]
\end{Lemma}

% block: negative-trace-space-intro page:38 type:ordinary-paragraph
{将此引理用于 $\gamma$ 等于迹算子 $\operatorname{tr}:H^1(\Omega)\to H^{1/2}(\partial\Omega)$ 的情形, 并将 $H^{-1/2}(\partial\Omega)$ 定义为 $H^{1/2}(\partial\Omega)$ 的对偶空间, 得到下述刻画.}

% block: trace-functional-representation page:38 type:proposition
\Needspace{6\baselineskip}
\begin{Proposition}\label{prop:trace-functional-representation}
{若 $L:H^1(\Omega)\to\mathbb R$ 是在 $\mathring H^1(\Omega)$ 上为零的有界线性泛函, 则存在唯一 $g\in H^{-1/2}(\partial\Omega)$, 使得}
\begin{equation}\label{eq:trace-functional-representation}
Lu=\langle g,\operatorname{tr}u\rangle,\qquad u\in H^1(\Omega).
\end{equation}
{此外, $\|g\|_{H^{-1/2}(\partial\Omega)}=\|L\|_{H^1(\Omega)'}$.}
\end{Proposition}

% block: trace-functional-proof page:38 type:proof
\begin{Proof}
{满足 \eqref{eq:trace-functional-representation} 的唯一 $g$ 的存在性直接来自该引理. 由 \eqref{eq:trace-functional-representation} 及 \eqref{eq:trace-quotient-norm}, 有}
\[
|Lu|\leq\|g\|_{H^{-1/2}(\partial\Omega)}\|\operatorname{tr}u\|_{H^{1/2}(\partial\Omega)}\leq\|g\|_{H^{-1/2}(\partial\Omega)}\|u\|_{H^1(\Omega)},
\]
{因此 $\|L\|\leq\|g\|_{H^{-1/2}(\partial\Omega)}$. 为证明反向不等式, 任取 $f\in H^{1/2}(\partial\Omega)$, 并选取 $u\in H^1(\Omega)$, 使 $\operatorname{tr}u=f$ 且 $\|u\|_1\leq\|f\|_{H^{1/2}(\partial\Omega)}$. 则}
\[
\langle g,f\rangle=Lu\leq\|L\|\|u\|_1\leq\|L\|\|f\|_{H^{1/2}(\partial\Omega)}.
\]
\end{Proof}

% block: normal-trace-functional page:39 type:ordinary-paragraph
{接下来应用此结论定义 $w\in H(\operatorname{div};\Omega)$ 的迹. 泛函}
\[
L_w(u)=\int u\operatorname{div}w\,dx+\int\operatorname{grad}u\cdot w\,dx,\qquad u\in H^1(\Omega),
\]
{是 $H^1$ 上的有界线性泛函, 且由分布导数的定义, 对 $u\in C_0^\infty(\Omega)$, $L_w(u)$ 为零, 从而由稠密性, 对 $u\in\mathring H^1(\Omega)$ 也为零. 因此可应用该命题得到唯一的 $\gamma_n w\in H^{-1/2}(\partial\Omega)$, 满足}
\begin{equation}\label{eq:normal-trace-green}
\int u\operatorname{div}w\,dx+\int\operatorname{grad}u\cdot w\,dx=\langle\gamma_n w,\operatorname{tr}u\rangle,\quad u\in H^1(\Omega),\ w\in H(\operatorname{div}).
\end{equation}

% block: normal-trace-norm-bound page:39 type:ordinary-paragraph
{此外,}
\[
\|\gamma_n w\|_{-1/2,\partial\Omega}=\sup_u\frac{\int u\operatorname{div}w\,dx+\int\operatorname{grad}u\cdot w\,dx}{\|u\|_{1,\Omega}}\leq\|w\|_{H(\operatorname{div})}.
\]

% block: normal-trace-classical-identification page:39 type:ordinary-paragraph
{为理解 $\gamma_n w$ 的含义, 将 \eqref{eq:normal-trace-green} 与分部积分公式 \eqref{eq:grad-div-classical-green} 比较. 固定 $w\in C^1(\overline\Omega;\mathbb R^3)$, 注意到作为 $u\in C^1(\overline\Omega)$ 的函数, \eqref{eq:grad-div-classical-green} 中的三个积分关于 $H^1$ 范数均有界. 由光滑函数在 $H^1$ 中稠密, \eqref{eq:grad-div-classical-green} 便对所有 $u\in H^1(\Omega)$ 成立. 然而 $\gamma_n w$ 由 \eqref{eq:normal-trace-green} 唯一确定, 因此有}
\[
\gamma_n w=w|_{\partial\Omega}\cdot n
\]
{其中 $w\in C^1(\overline\Omega;\mathbb R^3)$. 我们已得到下面的定理.}

% block: normal-trace-extension page:39 type:theorem
\Needspace{6\baselineskip}
\begin{Theorem}\label{thm:normal-trace-extension}
{$C^1(\overline\Omega;\mathbb R^3)$ 上的算子 $w\mapsto w|_{\partial\Omega}\cdot n$ 可延拓为有界线性算子 $\gamma_n:H(\operatorname{div})\to H^{-1/2}(\partial\Omega)$. 该延拓满足分部积分公式 \eqref{eq:normal-trace-green} 以及界 $\|\gamma_n w\|_{-1/2,\partial\Omega}\leq\|w\|_{H(\operatorname{div})}$.}
\end{Theorem}

% block: zero-normal-trace-domain page:39 type:ordinary-paragraph
{现在可以计算定义域为 $H^1$ 的梯度算子的伴随. 为此, 定义}
\[
\mathring H(\operatorname{div})=\{\,w\in H(\operatorname{div})\mid\gamma_n w=0\,\}.
\]

% block: gradient-adjoint page:39 type:theorem
\Needspace{6\baselineskip}
\begin{Theorem}\label{thm:gradient-adjoint-zero-normal}
{无界算子 $(\operatorname{grad},H^1)$ 的伴随为算子 $(-\operatorname{div},\mathring H(\operatorname{div}))$.}
\end{Theorem}

% block: gradient-adjoint-proof page:39 type:proof
\begin{Proof}
{向量场 $w\in L^2(\Omega;\mathbb R^n)$ 属于伴随的定义域, 当且仅当存在 $v\in L^2(\Omega)$, 使得}
\[
\int w\cdot\operatorname{grad}u\,dx=\int vu\,dx,\qquad u\in H^1(\Omega).
\]
{此时, 显然必须有 $w\in H(\operatorname{div})$ 且 $v=-\operatorname{div}w$ (考虑 $u\in C_0^\infty$ 即可). 因而 $w$ 属于伴随的定义域, 当且仅当 $w\in H(\operatorname{div})$ 且}
\[
\int w\cdot\operatorname{grad}u\,dx=-\int(\operatorname{div}w)u\,dx,\qquad u\in H^1(\Omega).
\]
{与 $\gamma_n w$ 的定义 \eqref{eq:normal-trace-green} 比较, 可见 $w$ 属于伴随的定义域当且仅当 $\gamma_n w=0$.}
\end{Proof}

% block: normal-trace-surjection-intro page:40 type:ordinary-paragraph
{接下来证明, 法向迹算子不仅将 $H(\operatorname{div})$ 映入 $H^{-1/2}(\partial\Omega)$, 而且映满该空间.}

% block: normal-trace-surjective page:40 type:theorem
\Needspace{6\baselineskip}
\begin{Theorem}\label{thm:normal-trace-surjective}
{算子 $\gamma_n:H(\operatorname{div})\to H^{-1/2}(\partial\Omega)$ 为满射.}
\end{Theorem}

% block: normal-trace-surjective-proof page:40 type:proof
\begin{Proof}
{给定 $g\in H^{-1/2}(\partial\Omega)$, 通过 Neumann 问题定义 $u\in H^1$:}
\[
\int\operatorname{grad}u\cdot\operatorname{grad}v\,dx+\int uv\,dx=\langle g,\operatorname{tr}v\rangle,\qquad v\in H^1.
\]
{则 $-\Delta u+u=0$. 取 $w=\operatorname{grad}u$. 则 $\operatorname{div}w=\Delta u=u$, 因此 $w\in H(\operatorname{div})$. 有}
\[
\int w\cdot\operatorname{grad}v\,dx+\int(\operatorname{div}w)v\,dx=\langle g,\operatorname{tr}v\rangle,\qquad v\in H^1,
\]
{这表明 $\gamma_n w=g$.}
\end{Proof}

% block: divergence-adjoint-intro page:40 type:ordinary-paragraph
{利用这一点, 可完全按照 Theorem~\ref{thm:gradient-adjoint-zero-normal} 证明中的论证, 证明下述结论, 从而确定散度算子的伴随.}

% block: divergence-adjoint page:40 type:theorem
\Needspace{6\baselineskip}
\begin{Theorem}\label{thm:divergence-adjoint-zero-trace}
{无界算子 $(\operatorname{div},H(\operatorname{div}))$ 的伴随为算子 $(-\operatorname{grad},\mathring H^1)$.}
\end{Theorem}

% block: reverse-grad-div-adjoints-intro page:40 type:ordinary-paragraph
{由于对闭的稠密定义算子有 $T^{**}=T$, 从 Theorems~\ref{thm:gradient-adjoint-zero-normal} 和 \ref{thm:divergence-adjoint-zero-trace} 可直接读出另外两个伴随.}

% block: reverse-grad-div-adjoints page:40 type:theorem
\Needspace{6\baselineskip}
\begin{Theorem}\label{thm:reverse-gradient-divergence-adjoints}
\[
\begin{gathered}
(\operatorname{div},\mathring H(\operatorname{div}))^*=(-\operatorname{grad},H^1),\\
(\operatorname{grad},\mathring H^1)^*=(-\operatorname{div},H(\operatorname{div})).
\end{gathered}
\]
\end{Theorem}

% block: grad-div-closed-ranges page:40 type:ordinary-paragraph
{下面验证这些算子的像空间全都是闭的. 对 $(\operatorname{div},H(\operatorname{div}))$, 这一点很明显, 因为实际上 $\operatorname{div}H(\operatorname{div})=L^2$: 将 $f\in L^2$ 写成某个 $u\in H^1$ 的 $\Delta u=\operatorname{div}\operatorname{grad}u$ 即可看出. 若 $\int f\,dx=0$, 则可选取 $u$ 满足 Neumann 条件 $\operatorname{grad}u\cdot n=0$, 从而 $f\in\operatorname{div}\mathring H(\operatorname{div})$. 因此像空间 $\operatorname{div}\mathring H(\operatorname{div})$ 在 $L^2$ 中的余维为 $1$, 所以由 Theorem~\ref{thm:finite-codimension-closed-range}, 它也是闭的. 闭值域定理进而表明这两个映射的伴随, 即 $(\operatorname{grad},\mathring H^1)$ 和 $(\operatorname{grad},H^1)$, 也都具有闭像空间.}

% block: curl-integration-by-parts-intro page:40 type:ordinary-paragraph
{最后转向 curl 算子, 假设区域 $\Omega\subset\mathbb R^3$. 上面对 grad 和 div 使用的大多数论证也适用于 curl, 但并非全部适用, 因为 curl 的核与像都不是有限余维的. 对 curl, 相应的分部积分公式为}
\begin{equation}\label{eq:curl-classical-green}
\begin{aligned}
\int v\cdot\operatorname{curl}u\,dx-\int\operatorname{curl}v\cdot u\,dx&=\int_{\partial\Omega}v|_{\partial\Omega}\times n\cdot u|_{\partial\Omega}\,ds,\\
&\hspace{-7em}u,v\in C^1(\overline\Omega;\mathbb R^3),
\end{aligned}
\end{equation}

% block: curl-distributional-definition pages:40,41 type:ordinary-paragraph
{其中 $n$ 为单位外法向量 (由于边界为 Lipschitz, 它几乎处处有定义). 空间 $H(\operatorname{curl})$ 定义为那些分布意义下的旋度也属于 $L^2(\Omega;\mathbb R^3)$ 的 $L^2(\Omega;\mathbb R^3)$ 向量场所成的空间. 因此, 向量场 $u\in L^2(\Omega;\mathbb R^3)$ 属于 $H(\operatorname{curl})$, 当且仅当存在向量场 $v\in L^2(\Omega;\mathbb R^3)$, 使得}
\begin{equation}\label{eq:distributional-curl}
\int_\Omega u\cdot\operatorname{curl}w\,dx=\int_\Omega v\cdot w\,dx,\qquad w\in C_0^\infty(\Omega;\mathbb R^3),
\end{equation}
{且此时 $v=\operatorname{curl}u$. 因此, 与上面相同,}
\[
(\operatorname{curl},C_0^\infty)^*=(\operatorname{curl},H(\operatorname{curl})).
\]
{作为伴随, 后一个算子是闭算子; 这意味着 $H(\operatorname{curl})$ 配备如下图范数后是 Hilbert 空间:}
\[
\|u\|_{H(\operatorname{curl})}^2=\|u\|_{L^2}^2+\|\operatorname{curl}u\|_{L^2}^2.
\]

% block: tangential-trace-functional page:41 type:ordinary-paragraph
{现在为 $v\in H(\operatorname{curl})$ 建立迹. 映射}
\[
u\mapsto\int v\cdot\operatorname{curl}u\,dx-\int\operatorname{curl}v\cdot u\,dx,\qquad u\in H^1(\Omega;\mathbb R^3),
\]
{是 $H^1(\Omega;\mathbb R^3)$ 上的有界线性泛函, 在 $C_0^\infty(\Omega;\mathbb R^3)$ 上为零, 因而由稠密性, 在 $\mathring H^1(\Omega;\mathbb R^3)$ 上也为零. 因此可应用 Proposition~\ref{prop:trace-functional-representation}, 推出存在唯一的 $\gamma_\tau v\in H^{-1/2}(\partial\Omega;\mathbb R^3)$, 使得}
\begin{equation}\label{eq:tangential-trace-green}
\begin{aligned}
\int v\cdot\operatorname{curl}u\,dx-\int\operatorname{curl}v\cdot u\,dx&=\langle\gamma_\tau v,\operatorname{tr}u\rangle,\\
u\in H^1(\Omega;\mathbb R^3),\quad v&\in H(\operatorname{curl}),
\end{aligned}
\end{equation}
{并得到范数界}
\begin{equation}\label{eq:tangential-trace-bound}
\|\gamma_\tau v\|_{-1/2,\partial\Omega}\leq\|v\|_{H(\operatorname{curl})}.
\end{equation}

% block: tangential-trace-identification page:41 type:ordinary-paragraph
{与 \eqref{eq:curl-classical-green} 比较可见, 我们已将\emph{切向迹算子} $\gamma_\tau v:=v|_{\partial\Omega}\times n$ 延拓为 $H(\operatorname{curl})$ 上的有界线性映射, 其像属于 $\partial\Omega$ 上向量场的 Sobolev 空间 $H^{-1/2}(\partial\Omega;\mathbb R^3)$. 因而得到了 $H(\operatorname{curl})$ 的相应于 Theorem~\ref{thm:normal-trace-extension} 的结论.}

% block: tangential-trace-extension page:41 type:theorem
\Needspace{6\baselineskip}
\begin{Theorem}\label{thm:tangential-trace-extension}
{$C^1(\overline\Omega;\mathbb R^3)$ 上的算子 $v\mapsto v|_{\partial\Omega}\times n$ 可延拓为有界线性算子 $\gamma_\tau:H(\operatorname{curl})\to H^{-1/2}(\partial\Omega;\mathbb R^3)$. 此延拓满足分部积分公式 \eqref{eq:tangential-trace-green} 和界 \eqref{eq:tangential-trace-bound}.}
\end{Theorem}

% block: tangential-trace-not-surjective page:41 type:ordinary-paragraph
{然而须注意, 迹算子 $\gamma_\tau$ 并不映满 $H^{-1/2}(\partial\Omega;\mathbb R^3)$, 即相应于 Theorem~\ref{thm:normal-trace-surjective} 的结论不成立. 首先, $\gamma_\tau w$ 总是切向向量场. 至少在 $v$ 和 $\partial\Omega$ 光滑时, 这一点很清楚, 而在边界仅为 Lipschitz 的情形, 如何精确定义这一点却很微妙. $\gamma_\tau$ 的像还受到其他限制. 详情参见 [\BookNumber{25}, \BookNumber{26}].}

% block: curl-adjoint-calculation pages:41,42 type:ordinary-paragraph
{最后计算 $(\operatorname{curl},H(\operatorname{curl}))$ 的伴随. 若 $w$ 属于伴随的定义域, 则存在 $v=(\operatorname{curl},H(\operatorname{curl}))^*w\in L^2(\Omega;\mathbb R^3)$, 满足}
\[
\int w\cdot\operatorname{curl}u\,dx=\int v\cdot u\,dx,\qquad u\in H(\operatorname{curl}).
\]
{这样的 $w$ 必须属于 $H(\operatorname{curl})$, 且 $\operatorname{curl}w=v$. 因此 $w$ 属于 $(\operatorname{curl},H(\operatorname{curl}))^*$ 的定义域, 当且仅当 $w\in H(\operatorname{curl})$ 且}
\[
\int w\cdot\operatorname{curl}u\,dx=\int\operatorname{curl}w\cdot u\,dx,\qquad u\in H(\operatorname{curl}).
\]
{可以证明 (用与 $H^1$ 情形类似的磨光论证), 光滑函数在 $H(\operatorname{curl})$ 中稠密. 因此只需要求上一个等式对所有 $u\in H^1(\Omega;\mathbb R^3)$ 成立. 而由 \eqref{eq:tangential-trace-green}, 这当且仅当 $\gamma_\tau w=0$ 时成立, 因而我们证明了 $(\operatorname{curl},H(\operatorname{curl}))$ 的伴随是 $(\operatorname{curl},\mathring H(\operatorname{curl}))$. 最后取伴随, 可得到反向关系. 这些结论概括在下面的定理中.}

% block: curl-adjoints page:42 type:theorem
\Needspace{6\baselineskip}
\begin{Theorem}\label{thm:curl-adjoints}
\[
\begin{gathered}
(\operatorname{curl},H(\operatorname{curl}))^*=(\operatorname{curl},\mathring H(\operatorname{curl})),\\
(\operatorname{curl},\mathring H(\operatorname{curl}))^*=(\operatorname{curl},H(\operatorname{curl})),
\end{gathered}
\]
{其中}
\[
\mathring H(\operatorname{curl}):=\{\,w\in H(\operatorname{curl})\mid\gamma_\tau w=0\,\}.
\]
\end{Theorem}

% block: curl-closed-range-forward-reference page:42 type:ordinary-paragraph
{在 Section~\ref{sec:l2-de-rham-complex} 中, 我们将证明像空间 $\operatorname{curl}H(\operatorname{curl})$ 是闭的.}

% block: unbounded-operators-summary-heading page:42 type:ordinary-paragraph
\begin{tcolorbox}[enhanced,breakable,colback=KuangNei,frame hidden,boxrule=0pt,sharpish corners,before skip=10pt,after skip=10pt]
\refstepcounter{booksummary}\label{box:unbounded-operators-summary}
\noindent\textbf{Box~\thebooksummary:} {Chapter~\ref{chap:unbounded-operators} 中 Hilbert 空间上算子的内容摘要.}\par

% block: summary-unbounded-domain page:42 type:ordinary-paragraph
{从 Hilbert 空间 $X$ 到另一个空间 $Y$ 的\emph{(无界) 算子}, 是将其定义域 $D(T)$ (它是 $X$ 的子空间) 映入 $Y$ 的线性算子. 定义域不必在 $X$ 中闭, 算子也不必有界, 但通常要求算子\emph{稠密定义}, 即 $D(T)$ 在 $X$ 中稠密.}

% block: summary-closed-operator page:42 type:ordinary-paragraph
{若算子的图在 $X\times Y$ 中闭, 或等价地, $D(T)$ 配备\emph{图范数}后完备, 则称该算子是\emph{闭的}. 处处有定义的算子是闭的当且仅当它有界.}

% block: summary-adjoint-operator page:42 type:ordinary-paragraph
{从 $X$ 到 $Y$ 的稠密定义算子 $T$ 产生从 $Y$ 到 $X$ 的伴随算子 $T^*$. 伴随 $T^*$ 总是闭的. 若 $T$ 是闭的, 则 $T^*$ 也稠密定义, 并且此时 $T^{**}=T$.}

% block: summary-closed-range page:42 type:ordinary-paragraph
{闭的稠密定义线性算子 $X\to Y$ 的像空间不一定是闭子空间, 但 $T$ 的像空间在 $Y$ 中闭当且仅当 $T^*$ 的像空间在 $X$ 中闭. 若 $\dim Y/\mathcal R(T)<\infty$, 则 $T$ 具有闭像空间. 若包含映射 $D(T)\subset X$ 是紧的, 则 $T$ 具有闭像空间.}

% block: summary-vector-operators page:42 type:ordinary-paragraph
{在 $\mathbb R^3$ 中具有 Lipschitz 边界的有界区域上, 算子 grad, curl 和 div 可视为 $L^2$ (标量或向量) 上的无界算子, 定义域分别为 $H^1$, $H(\operatorname{curl})$ 和 $H(\operatorname{div})$. 它们稠密定义, 闭, 且具有闭像空间. 它们的伴随分别为 $-\operatorname{div}$, $\operatorname{curl}$ 和 $-\operatorname{grad}$, 定义域分别为 $\mathring H(\operatorname{div})$, $\mathring H(\operatorname{curl})$ 和 $\mathring H^1$.}
\end{tcolorbox}
